\section{Local projections and amplitude levels}\label{sec:geometry}

For the remainder of the proof suppose \(2<s<d\).
Partition \(\T^d\) into half-open dyadic cubes \(I\) of side
\(h\le1/8\). For each \(I\), enlarge its Euclidean lift by a
factor of three about its center and denote the resulting cube
by \(J_I\), with volume \(v=(3h)^d\). Its image \(J_I^{\rm phys}\)
on the torus is Borel, and the quotient restricted to the lift
is injective and measure preserving. Integrals and restrictions
to \(J_I\) refer to this lift with the periodized density. The
physical supports have overlap at most \(5^d\). Indeed, the center
of any support containing a point lies within circular distance
\(3h/2\) in each coordinate on a grid of spacing \(h\). Two
intersecting supports have centers within circular distance \(3h\)
in every coordinate. Thus the graph joining intersecting physical
supports has maximum degree at most \(7^d-1\).

Use real polynomials of total degree at most \(r=\lceil s\rceil\).
Let \(Q_I\) be the ordinary \(L^2\) projection of \(f\) on \(J_I\)
and \(e_I=\|f-Q_I\|_{2,J_I}\). The kernel on the volume-one cube
\(\Omega=[0,1)^d\) is \(K_\Omega(x,y)=\sum_i\phi_i(x)\phi_i(y)\)
for a real orthonormal basis. It is uniformly bounded and
reproduces the polynomial space. Rescaling the kernel shows that
\(Q_I\) has mean \(b_I=P_I/v\) and supremum norm at most \(C_V b_I\),
where \(P_I=\int_{J_I} f\). Write \(\mu_I=\int_I f\).
Under the affine map \(T_I\) from \(\Omega\) into the torus
through \(J_I\), the central cube \(W=[1/3,2/3)^d\)
maps to \(I\). The volume Jacobian of \(T_I\) is \(v\),
and \(vb_I=P_I\). If \(P_I>0\), set
\[
 w_I=f\circ T_I/b_I,\qquad q_I=Q_I\circ T_I/b_I,\qquad
 E_I=\|w_I-q_I\|_{2,\Omega}.
\]
Then \(w_I\) is a probability density, \(q_I\) is its ordinary
projection, \(\mu_I=P_I\int_W w_I\), and \(P_I E_I=\sqrt v\,e_I\).
If \(P_I=0\) the density and projection vanish almost everywhere on
\(J_I\), and products \(P_I E_I\) are defined to be zero.

We record the projection errors needed below. For every dyadic
\(h\le1/8\) and each dyadic \(\ell\in[h,1]\), let \(P_\ell\) be the
ordinary piecewise projection on disjoint cubes of side \(\ell\).
With gradients taken separately in the cubes,
\begin{equation}\label{eq:projection}
 \sum_I\sqrt v\,e_I\le C h^s,\qquad
 \|\nabla(f-P_\ell)\|_{2,\mathrm{piecewise}}\le C\ell^{s-1}.
\end{equation}
To see this, at side \(\ell\) split \(f\) at frequency of order
\(\ell^{-1}\). The tail and its gradient are \(O(\ell^s)\) and
\(O(\ell^{s-1})\) in \(L^2\). Indeed \eqref{eq:tail}
gives squared shell bounds \(C2^{-2js}\), and the
squared gradient tail is bounded by
\(C\sum_{2^j\gtrsim \ell^{-1}} 2^{-2j(s-1)}\).
In particular \(f\in W^{1,2}(\T^d)\). For the smooth
part \(g\), on any cube of side \(\ell\) choose the degree-\(r\)
polynomial \(P\) for which the derivatives of \(g-P\) through
order \(r\) have mean zero. Its coefficients are obtained in
descending order of total degree using the nonzero constant
derivative of the matching monomial. Poincaré's inequality on
each derivative with mean zero, iterated, gives
\[
 \|D^j(g-P)\|_2\le C\ell^{r+1-j}\|D^{r+1}g\|_2,\qquad j=0,1 ,
\]
on that cube, with the norms summing over derivatives of each order.
If \(\Pi\) is the local orthogonal projection, then
\(g-\Pi g=(g-P)+\Pi(P-g)\). Apply \(L^2\) contraction and
\(\|\nabla \Pi(P-g)\|_2\le C\ell^{-1}\|\Pi(P-g)\|_2\)
to retain the estimates for \(j=0,1\). The inverse inequality
follows by rescaling and finite-dimensional norm equivalence.
It also bounds the gradient of the projected tail. Over the torus,
\(\|D^{r+1}g\|_2\le C\ell^{s-r-1}\) by summing shells, since
\(r+1>s\). The argument works on the bounded-overlap family
\(J_I\) at side \(3h\) as well, giving
\(\sum_I e_I^2\le C h^{2s}\). Cauchy--Schwarz and
\(\sum_I v=3^d\) give the first bound in \eqref{eq:projection}.
Every estimate holds uniformly over the class on this mesh range.

\begin{lemma}\label{lem:sublevel}
For fixed degree bound \(r\), there are \(\beta>0\) and \(C<\infty\)
such that any real polynomial \(q\) of degree at most \(r\) on
\(\Omega\) with \(\|q\|_1\ge1\) satisfies
\[
 |\{x\in\Omega:|q(x)|\le z\}|\le C z^\beta
 \quad(z>0),\qquad \|\log^-|q|\|_{2,\Omega}\le C .
\]
\end{lemma}
\begin{proof}
In one dimension, if the sublevel set has length \(\lambda>0\),
choose \(r+1\) points from it with successive pairwise distances at
least \(\lambda/(4(r+1))\) by excluding intervals around the
points already chosen. Lagrange interpolation yields
\(\|q\|_\infty\le C z\lambda^{-r}\). Constant polynomials satisfy
the claim directly. For induction on dimension, write
\(q(x_1,x')=\sum_{i=0}^r x_1^i q_i(x')\).
Some coefficient has \(L^1\) norm at least \(1/(r+1)\). Outside
a set of measure \(O(z^{\beta'/2})\), its magnitude is at least
\(\sqrt z\) for \(0<z<1\), by the induction bound with exponent
\(\beta'\), scaled by its norm. A slice with such a coefficient
has supremum at least \(c\sqrt z\) by finite-dimensional norm
equivalence. The interval bound and Fubini prove the induction.
The zero set is null, and the tail of \(\log^-|q|\) is bounded
by \(C e^{-\beta t}\) for \(t\ge0\). Integrating the squared
tail formula gives the second bound.
\end{proof}

Put \(\psi(z)=z\log|z|\) for \(z\ne0\) and \(\psi(0)=0\).
Our deterministic surrogate on cell \(I\) is
\begin{equation}\label{eq:surrogate}
 F_I=P_I\int_W\psi(q_I)+\mu_I\log b_I
\end{equation}
for positive mass, and zero otherwise. For \(u\ge0\) and \(z\ne0\)
the divided difference of \(\psi\) in absolute value is at most
\[
 C+\log^+(u+|z|)+\log^-|z|.
\]
For \(z>0\) it is \(\log z + x\log x/(x-1)\) when \(x=u/z<1\),
and \(\log u+\log x/(x-1)\) when \(x>1\). The logarithmic ratios
here belong to \([0,1]\), including the endpoint limits. For \(z<0\)
the quotient is \(\log(u+|z|)+x\log x+(1-x)\log(1-x)\),
where \(x=u/(u+|z|)\). The last two terms have bounded sum.
As \(w_I\) has mass one and \(q_I\) has bounded \(L^1\) norm and
mean one, \((\log^+ t)^2\le C t\) and Lemma~\ref{lem:sublevel}
bound the divided difference in \(L^2(W)\) uniformly. Thus
Cauchy--Schwarz and the change of variables to \(W\) yield
\begin{equation}\label{eq:transfer}
 \sum_I\left|\int_I f\log f-F_I\right|
 \le C\sum_I P_I E_I\le C h^s .
\end{equation}
The scale term \(\mu_I\log b_I\) is common to the integral and
the surrogate, since it is formed with the mass of \(w_I\) on \(W\).

For \(P_I>0\) set \(\nu_I=\one_{\{\inf_W|q_I|<1/2\}}\),
and take \(\nu_I=0\) at zero mass.
This indicator records when the infimum of the absolute projected value
on \(I\) is below half the mean on \(J_I\).
Set \(\tau=h^s\) and \(\omega=2/(2s-1)\).
\begin{lemma}\label{lem:bands}
For \(B\ge\tau\), let \(\mathcal V_B\) contain the cells with
\(\nu_I=1\) and \(B/2<b_I\le B\). At \(\xi=B/\tau\),
\[
 h^d\#\mathcal V_B\le
 C\left(\xi^{-\omega}+h+
         B^{-1}\sum_{B/2<b_I\le B}\sqrt v\,e_I\right).
\]
\end{lemma}
\begin{proof}
Take \(U=f\wedge (C_0B)\) with \(C_0\ge \max(1,2C_V)\).
Let \(\Phi(y)=y^2\) on \(0\le y\le1\) and \(\Phi(y)=2y-1\)
on \(y>1\). On each \(J_I\),
\[
 \|U-\overline U_{J_I}\|_2^2
 \le C B^2\sum_{j=1}^d\int_{J_I}\Phi(h|\partial_j U|/B).
\]
On a coordinate slice, a pairwise difference is bounded by both
\(C_0 B\) and the integral of the absolute derivative on the whole
slice. Its square is bounded by \(CB^2\Phi\) of the latter integral
divided by \(B\). Use Jensen along the slice, whose length is \(3h\),
and average the pairwise differences. In the product cube, telescope
in coordinates and sum their squares, or use the coordinate
averaging contractions, to obtain the displayed inequality.

On each macro cube of side \(\ell\), fix a slice in direction \(j\)
and write \(g=\partial_j f\), \(k=\partial_j P_\ell\), \(M=B/h\).
The scalar bound
\[
 \Phi(|g|/M)\le
 C\bigl(\cl(k/M,[-1,1])\,g/M+(g-k)^2/M^2\bigr)
\]
holds for all real \(g,k\). For \(|k|\le M\), completing the
square controls \(g^2/M^2\). For \(k/M\ge1\), the right side
without the constant is \(y+(y-k/M)^2\) with \(y=g/M\).
For \(y\ge1/2\) use the linear term, for \(0\le y<1/2\)
use the squared term, and for \(y<0\) use
\(-|y|+(|y|+1)^2\). Reversing both signs covers \(k/M\le-1\).
Apply the inequality on \(\{f<C_0 B\}\), where
\(\partial_j U=g\). Elsewhere \(\partial_j U=0\) almost
everywhere by the Sobolev chain rule for a truncation.
The integral of \(\cl(k/M,[-1,1])\,\partial_j U/M\) on the
slice is at most \(C B/M\). Indeed the clipped polynomial
has bounded variation uniformly, from its finitely many
monotonicity intervals, and integration by parts also bounds the
two endpoint contributions. Sobolev slicing justifies the
one-dimensional integrals. The cross-sectional volume of each
macro cube is \(\ell^{d-1}\) and there are \(\ell^{-d}\) cubes,
so the linear term contributes \(O(h/\ell)\). The squared
term contributes at most \(M^{-2}\|\nabla(f-P_\ell)\|_2^2\)
up to a constant. Hence \eqref{eq:projection} gives
\begin{equation}\label{eq:huber}
 \sum_j\int_{\T^d}\Phi(h|\partial_j U|/B)
 \le C\left(h/\ell+h^2\ell^{2s-2}/B^2\right).
\end{equation}

In the band under study, \(Q_I\le C_VB\), so
\(\|f-U\|_{2,J_I}\le e_I\) and \(\|U-Q_I\|_{2,J_I}\le2e_I\).
The centered functions \(U-\overline U_{J_I}\) and \(Q_I-b_I\)
are bounded by \(CB\). Comparing squares and using
Cauchy--Schwarz gives
\(\|Q_I-b_I\|_{2,J_I}^2\le
\|U-\overline U_{J_I}\|_2^2+CB\sqrt v\,e_I\).
For \(\nu_I=1\) the supremum of \(|Q_I-b_I|\) is at least
\(b_I/2\), and by norm equivalence its squared \(L^2(J_I)\)
norm is at least \(cB^2 v\). Sum, use the bounded overlap
and \eqref{eq:huber}, then divide by \(B^2\). This proves the
claimed estimate with \(h/\ell+h^2\ell^{2s-2}/B^2\)
in place of \(\xi^{-\omega}+h\). When \(h\xi^\omega\le1\),
take a dyadic \(\ell\) within a factor of two of \(h\xi^\omega\)
in \([h,1]\). The two terms are \(O(\xi^{-\omega})\).
When \(h\xi^\omega>1\), use \(\ell=1\) and
\(h^{2-2s}\xi^{-2}\le h\).
\end{proof}
